\documentclass{article}

\usepackage{makecell}

\title{Ch. 15 HW \#2 -- Solutions}
\author{Principles of Microeconomics}
\date{\today}

\begin{document}

\maketitle

\begin{enumerate}

\setcounter{enumi}{3}
\item (d) -- Profit maximizing firms produce the quantity at which marginal revenue (price) equals marginal cost, so the marginal cost curve determines the quantity supplied for a given price. In the short run, firms' fixed costs are lost regardless of whether they produce, so they will produce as long as their revenue is greater than their variable costs. 

\item (a) -- Profit maximizing firms produce the quantity at which marginal cost equals price. That means if marginal cost is between average variable cost and average total cost, price is between average variable cost and average total cost. Therefore, the firm makes enough revenue to offset their variable costs but not their total costs. In the short run, they will keep producing to minimize their losses, but in the long run, they will exit because they're unprofitable.

\section*{Quick Quiz}

\section*{Problems \& Applications}

\setcounter{enumi}{1}
\item Leah's variable cost per day is \[\$280 - \$30 = \$250.\] Therefore, Leah's average variable cost is \[\$250 / 10 = \$25.\] That's less than her average revenue of \$27, so Leah produces in the short run.

In the long run, Leah's average total cost is \[\$280 / 10 = \$28.\] That's greater than her average revenue, \$27, so Leah exits in the long run.

\setcounter{enumi}{3}
\item

	\begin{enumerate}
	
	\item \phantom{x}
	
		\begin{center}
		\begin{tabular}{| c | c | c | c | c | c | c | c |}
		\hline
		Q & FC & VC & AFC & AVC & ATC & MC & Profit \\
		\hline
		0 & 100 & 0 & -- & -- & -- & -- & -100 \\
		\hline
		1 & 100 & 50 & 100 & 50 & 150 & 50 & -100 \\
		\hline
		2 & 100 & 70 & 50 & 35 & 85 & 20 & -70\\
		\hline
		3 & 100 & 90 & 33.33 & 30 & 63.33 & 20 & -40 \\
		\hline
		4 & 100 & 140 & 25 & 35 & 60 & 50 & -40 \\
		\hline
		5 & 100 & 200 & 20 & 40 & 60 & 60 & -50\\
		\hline
		6 & 100 & 360 & 16.67 & 60 & 76.67 & 160 & -160\\
		\hline
		\end{tabular}
		\end{center}
		
	\item See the profit column above. No, shutting down is not a wise decision because their losses would be greater from shutting down than they would be from producing 4 cases.
	
	\item The firm's losses are \$100. It's not the best decision because if the firm produces 4 cases -- the other quantity at which marginal cost equals marginal revenue -- its losses are smaller, \$40. That's because the firm's marginal costs decline at first, so the firm's marginal costs are actually less than its marginal revenue for cases 2 \& 3. 
	
	\end{enumerate}

\setcounter{enumi}{5}	
\item In a competitive market, all of a firm's output gets sold at the market price. That means marginal revenue is equal to average revenue which is equal to the market price of the good, Therefore, the average revenue and the price of the good both equal \$10. Since the price is \$10, that means the firm sold $\$500 / \$10 = \$50$

\item

	\begin{enumerate}
	
	\item Price equals average revenue which equals \$10. Therefore, profit is \[\pi = (P - ATC)*Q = (10 - 8)*100 = 200.\]
	
	\item Since the firm is profit maximizing, marginal cost equals marginal revenue which equals the price of the good, \$10.
	
	\item Average variable cost is 
	\[ AVC = ATC - AFC = ATC - FC/Q = 8 - 200/100 = \$6.\]
	
	\item Since marginal cost is greater than average total cost, average total cost is greater than its minimum. That means the firm's efficient scale is less than 100 units.
	
	\end{enumerate}

\end{enumerate}

\end{document}